\section{Keamanan Komputasi dan Side-Channel Attack}
Sebuah algoritma dikatakan \textbf{aman secara komputasi (\textit{computationally secure})} jika:
1. Persamaan matematikanya terlalu kompleks untuk dipecahkan secara analitik.
2. Biaya untuk memecahkan cipherteks melebihi nilai informasi di dalamnya.
3. Waktu yang dibutuhkan untuk memecahkan cipherteks melebihi masa berlaku kerahasiaan informasi tersebut.

\subsection{Side-Channel Attack}
Serangan ini tidak menyerang algoritma matematika, melainkan memanfaatkan kebocoran fisik saat algoritma dijalankan pada perangkat:
\begin{itemize}
    \item \textbf{Timing Attack}: Menganalisis perbedaan waktu eksekusi.
    \item \textbf{Power Analysis}: Menganalisis konsumsi daya (SPA/DPA).
    \item \textbf{Electromagnetic Attack}: Menganalisis radiasi elektromagnetik.
    \item \textbf{Fault Injection}: Sengaja membuat kesalahan pada perangkat untuk menganalisis output yang salah (DFA).
\end{itemize}





